\section*{Sets and Indices}

\[
C = \{\text{corn},\text{wheat},\text{soybeans},\text{sorghum}\}
\]

No additional indexing sets are used in the implemented model.

\section*{Parameters}

For each crop \(c \in C\), let \(p_c\) denote profit per acre:
\[
p_{\text{corn}}=1500,\quad
p_{\text{wheat}}=1200,\quad
p_{\text{soybeans}}=1800,\quad
p_{\text{sorghum}}=1600.
\]

Other scalar parameters are:
\[
A = 100
\qquad\text{(total farm area, code: \texttt{total\_area})}
\]
\[
r^{cw} = 2
\qquad\text{(corn--wheat ratio, code: \texttt{corn\_wheat\_ratio})}
\]
\[
r^{ss} = 0.5
\qquad\text{(soybeans--sorghum ratio, code: \texttt{soybeans\_sorghum\_ratio})}
\]
\[
r^{ws} = 3
\qquad\text{(wheat--sorghum ratio, code: \texttt{wheat\_sorghum\_ratio})}
\]
\[
L^{\min}_{\text{corn}} = 20
\qquad\text{(minimum corn acres, code: \texttt{min\_corn\_acres})}.
\]

\section*{Decision Variables}

Continuous acreage variables:
\[
x_{\text{corn}} \ge 0 \quad \text{(code: \texttt{corn})}
\]
\[
x_{\text{wheat}} \ge 0 \quad \text{(code: \texttt{wheat})}
\]
\[
x_{\text{soybeans}} \ge 0 \quad \text{(code: \texttt{soybeans})}
\]
\[
x_{\text{sorghum}} \ge 0 \quad \text{(code: \texttt{sorghum})}.
\]

Binary activation variables for the mutual-exclusion logic:
\[
y_{\text{corn}} \in \{0,1\}
\quad \text{(code: \texttt{y\_corn})}
\]
\[
y_{\text{soybeans}} \in \{0,1\}
\quad \text{(code: \texttt{y\_soybeans})}.
\]

\section*{Objective}

Maximize total profit:
\[
\max \;
p_{\text{corn}}x_{\text{corn}}
+ p_{\text{wheat}}x_{\text{wheat}}
+ p_{\text{soybeans}}x_{\text{soybeans}}
+ p_{\text{sorghum}}x_{\text{sorghum}}.
\]

\section*{Constraints}

Total available land:
\[
x_{\text{corn}} + x_{\text{wheat}} + x_{\text{soybeans}} + x_{\text{sorghum}} \le A.
\]

Corn-to-wheat acreage ratio:
\[
x_{\text{corn}} \ge r^{cw} x_{\text{wheat}}.
\]

Soybeans-to-sorghum acreage ratio:
\[
x_{\text{soybeans}} \ge r^{ss} x_{\text{sorghum}}.
\]

Wheat-to-sorghum equality:
\[
x_{\text{wheat}} = r^{ws} x_{\text{sorghum}}.
\]

Minimum corn commitment:
\[
x_{\text{corn}} \ge L^{\min}_{\text{corn}}.
\]

Binary activation for corn:
\[
x_{\text{corn}} \le A\, y_{\text{corn}}.
\]

Binary activation for soybeans:
\[
x_{\text{soybeans}} \le A\, y_{\text{soybeans}}.
\]

Mutual exclusion of corn and soybeans:
\[
y_{\text{corn}} + y_{\text{soybeans}} \le 1.
\]

\section*{Notes on Implementation}

The code solves the above formulation as a mixed-integer linear program using PuLP-compatible modeling objects and a Gurobi backend (\texttt{GUROBI\_CMD}).

The mutual exclusion is implemented exactly through the binary variables and big-\(M\) activation constraints with \(M=A=100\). Because of
\[
x_{\text{corn}} \ge 20
\quad\text{and}\quad
x_{\text{corn}} \le A y_{\text{corn}},
\]
the model forces \(y_{\text{corn}}=1\). Then
\[
y_{\text{corn}} + y_{\text{soybeans}} \le 1
\]
implies \(y_{\text{soybeans}}=0\), and hence
\[
x_{\text{soybeans}} \le 0.
\]
So in the implemented model, soybeans are completely excluded whenever the minimum-corn requirement is present.

No integer acreage restrictions are imposed; acreages are continuous nonnegative variables. The wheat--sorghum relationship is modeled as an equality, not an inequality.
